\section{Anhang: COLMAP-Vorstudie}

Die COLMAP-Voruntersuchung wurde als separater Benchmark durchgeführt und ist
nicht mit der späteren GPU-Gesamtpipeline gleichzusetzen. Sie bewertet nur die
Sparse-SfM-Stufe auf dem Alurohr-Testdatensatz. Alle fünf dokumentierten Läufe
registrierten 240 von 240 Bildern. Die Auswertung ist in
\texttt{docs/COLMAP\_Tests/} archiviert.

\begin{table}[H]
\centering
\small
\begin{tabularx}{\textwidth}{Xrrrrr}
\toprule
Konfiguration & Punkte & Punkte/Bild & Tracks & Reprojektion & Zeit (s)\\
\midrule
Plain-SIFT 4096 & 139449 & 581 & 6.69 & 0.693 px & 876\\
Plain-SIFT 4096 + Guided & 142208 & 593 & 7.44 & 0.790 px & 1678\\
Plain-SIFT 8192 & 140594 & 586 & 6.71 & 0.698 px & 925\\
CRF 18 & 94616 & 394 & 6.64 & 0.728 px & 677\\
CRF 28 & 93703 & 390 & 6.43 & 0.761 px & 725\\
\bottomrule
\end{tabularx}
\caption{Optionale COLMAP-SfM-Voruntersuchung. Die Läufe sind ein separater
Benchmark und kein vollständiger STS-/SuGaR-Vergleich.}
\label{tab:colmap-vorstudie}
\end{table}

Die Ergebnisse begründen den Projektarbeitsstandard Plain-SIFT mit 4096
Merkmalen, Sequential-Overlap 15 und deaktiviertem Guided Matching. Guided
Matching brachte nur einen kleinen Punktezuwachs, erhöhte jedoch die Laufzeit
stark und verschlechterte den mittleren Reprojektionsfehler. Die vollständigen
Rohberichte liegen in den Dateien unter
\texttt{docs/COLMAP\_Tests/}.
